\begin{problem}[第36届全国中学生物理竞赛决赛][超导重力仪与地下水重力变化]{113}
    用超导重力仪对地球表面重力加速度的微小变化及其分布进行长期高精度观测，可获得地幔运动、质量分布、固体潮汐、地下水和矿产分布等地球内部结构信息，以及海洋潮汐、气候变化等外部数据。超导重力仪灵敏度高（可达地表重力加速度的 $10^{-9}$）、稳定性强，是目前地球科学领域广泛采用的重力观测装置。随着分布于全球的超导重力仪数目的增加、观测网络的建立和数据共享，为地球科学的发展、万有引力规律的精确检验等等提供了新的机遇。试就有关地下水分布引起的重力变化、超导重力仪原理等问题给出回答。
    \begin{subproblem}
        假设地球可视为一个半径为 $6370\,\mathrm{km}$ 的均匀球体，其表面重力加速度 $g_{0}=9.80\,\mathrm{m\cdot s^{-2}}$（不考虑地球自转）。现某处地下有一个直径为 $20\,\mathrm{km}$ 充满水的球形地下湖，其球心离地面 $15\,\mathrm{km}$。求地下湖正上方地表处重力加速度的微小变化 $\Delta g$。已知万有引力常量 $G=6.67\times10^{-11}\,\mathrm{N\cdot m^{2}\cdot kg^{-2}}$，水的密度 $\rho_{\text{水}}=1.0\times10^{3}\,\mathrm{kg\cdot m^{-3}}$。
    \end{subproblem}
    \begin{subproblem}
        超导重力仪的核心部件是悬浮在磁场中、由金属铌（Nb）制成的超导球壳，为计算方便起见，这里简化为超导细环（见图 a）。悬浮磁场由准亥姆霍兹（Helmholtz）线圈提供，它由匝数分别为 $\alpha N$ 和 $N$ 的上、下两个同轴细超导线圈（$\alpha < 1$）串联而成，两线圈的半径及其中心之间的距离均为 $R$。超导重力仪调试前，亥姆霍兹线圈和超导细环中的电流均为零，超导细环处在下线圈平面内（$z = 0$），且与线圈同轴。调试开始后，外电源对准亥姆霍兹线圈缓慢加上电流 $i_{0}$，在此过程中超导细环内会产生感应电流。调节 $i_{0}$ 的大小，使超导细环正好悬浮在上线圈平面内（$z = R$），如图a所示。由于超导线圈和超导细环的电阻均为零（其所在处的温度为 $4.2\mathrm{K}$），超导环被稳定悬浮后，移去提供超导线圈电流的电源，超导线圈中的电流 $i_0$ 以及由此产生的磁场均能长期保持稳定。已知超导细环的质量为 $m$、直径为 $D(D \ll 2R)$、自感系数为 $L$；在上线圈平面附近，由准亥姆霍兹线圈产生的磁感应强度 $B$ 的 $z$ 方向分量和径向分量可分别近似表示为：
        \[
        \left\{ \begin{array}{l}
        B_{z} = B_{0}[1 - \beta(z - R)] \\
        B_{r} = \frac{1}{2} \beta B_{0} r
        \end{array} \right.
        \]
        其中 $B_{0}$ 为上线圈中心处的磁感应强度，$\beta$ 为径向系数，$r$ 是到轴线的距离。利用上述线圈和超导细环的已知参数、以及线圈中的外加电流 $i_0$，求
        \begin{subsubproblem}
            $B_{0}$ 和 $\beta$ 的表达式；
        \end{subsubproblem}
        \begin{subsubproblem}
            超导细环在平衡处（$z = R$）感应电流 $I_{0}$ 的表达式；
        \end{subsubproblem}
        \begin{subsubproblem}
            超导重力仪放置地的重力加速度 $g_{0}$ 的表达式。
        \end{subsubproblem}
    \end{subproblem}
    \begin{subproblem}
        为了精确测量重力加速度的微小变化，在超导重力仪中采用了交流电桥平衡法。为了理解其原理并便于计算，将超导细环视为面积为 $A$ 的薄“平板”，并在超导细环的上下对称位置，即 $R + d/2$ 和 $R - d/2$ 处，分别固定两个形状相同、面积也为 $A$ 的金属平板，构成上极板与超导“平板”、下极板与超导“平板”两个电容器 $C_{1}$ 和 $C_{2}$，如图 b 所示。再将 $C_{1}$ 和 $C_{2}$ 与两组完全相同的电感 $L_{0}$ 和电容 $C_{0}$ 连接成电桥，并接在交流电源（频率远高于重力加速度的变化的快慢）上，如图 c 所示。假设某时刻超导重力仪所在地的重力加速度为 $g$，超导细环（超导“平板”）处在其平衡位置 $z = R$ 处，交流电桥处于平衡，这时 $C_{1}$ 和 $C_{2}$ 上加有相同的直流电压 $V_{C}$；当重力加速度改变 $\Delta g$，导致电桥失衡，将 $C_{1}$ 上的直流电压增加 $\Delta V$，同时在 $C_{2}$ 上减少 $\Delta V$，电桥又重新达到平衡。试给出 $\Delta g$ 与 $\Delta V$ 的关系式（不计电容器的边缘效应）。
    \end{subproblem}
\end{problem}